% Conclusion and Future Work (about 6 pages). Written in Stage 7. No new numbers: every figure repeats one stated
% in Chapters 5 and 7, with its row in evidence_registry.csv.
\chapter{Conclusion and Future Work}\label{ch09}

\section{Summary of Findings}\label{sec:ch09-summary}

This thesis built and evaluated a system that turns mixed tables and text into a knowledge graph and answers
natural-language questions over it. A language model proposes the plan for the tables, the kinds of facts in the
text, the facts and mentions themselves, and the identity of near misses; code validates every proposal, verifies
every quotation, and decides wherever a rule can. A tool-using agent answers questions by calling read-only tools one
at a time, and code checks every call, every citation and every outcome.

The central finding concerns what the graph should be. Four representations were built and measured in turn. Storing
extracted facts as edges between shared nodes leaked facts from one thing to another (36 of 100 audited paths false
on the furniture data). Storing each fact as a node removed the leak but showed how little of a text the stored facts
cover, and a richer layered model of facts did not improve answers over records combined with vector retrieval. On
the final builds, stored facts are almost always right (96--99\,\% content precision) but cover only 18--58\,\% of the
claims a careful reader finds, mostly because a proposed schema cannot anticipate what a corpus says. The final design
therefore uses the graph as an \emph{anchor graph}: an index that locates the records, individuals and concepts a
question is about and leads from them to the passages that concern them, with identity stored as reversible,
evidenced links and a rule that a shared node finds passages but never proves a fact. Judged as an index, by criteria
fixed in advance, the anchor graph passed every hard gate on four datasets.

The query engine answered 425 of 440 unseen questions right by meaning (96.6\,\%; 421 by exact match) on those four
datasets, far more than retrieval-only readers on the same questions, at about 1.6 cents and 10 seconds per question.
Its remaining errors lie mostly in how the model reads: passages that state a missing answer in other words are found
and then left out.

\section{Answers to the Research Questions}\label{sec:ch09-answers}

\paragraph{RQ1 --- Can an LLM-proposed, code-verified pipeline build a correct, entity-centric, provenance-aware graph
from heterogeneous data without domain-specific code?} Yes, on the four datasets measured and within their form. The
same code built all four graphs; every item made from text carries a verified quotation; no two records or
individuals were merged wrongly; every judged link from a mention to a record was right (36/36, 63/63, 77/77, 48/48);
and walks from a record or individual reached only the chunks that concern it, with one exception on documents that
describe several events. These gates were reached by development on three datasets and held on the fourth, unseen one
without change; a fifth dataset that failed one gate shows that they are not guaranteed (\cref{ch05}).

\paragraph{RQ2 --- What are the trade-offs between storing extracted facts and anchoring evidence?} Stored facts are
precise but incomplete, can leak across shared nodes, and depend on a schema that cannot anticipate a new corpus; a
graph built on them did not answer better than records with vector retrieval. An anchor graph gives up exact answers
from stored facts --- which the recall measurements show were not available --- and gains identity, links and
provenance that can be checked. Navigating through the stored facts added no evidence and added unwitnessed
connections. The anchor graph is not a better ranker of evidence than vector search; its value is in saying correctly
what a passage is about and which record a name refers to (\cref{ch04,ch05}).

\paragraph{RQ3 --- How accurately can a tool-using agent answer questions over tables and text, with verifiable
citations?} With 96.6\,\% accuracy by meaning on 440 unseen questions, and far more than retrieval-only readers on paired
sets (for example 63 against at most 31 of 68 development questions on the NHTSA data), most clearly on questions
about records, joins and counts. Every answer cites identifiers that tools returned, and code checks them. The
accuracy belongs to the engine versions measured; the final engine was measured on answers only on small sets, and
the share of accuracy the graph itself provides was not isolated (\cref{ch07,ch08}).

\paragraph{RQ4 --- How do retrieval, tool design and reasoning budget affect accuracy, cost and latency?} Fused card
lists find start nodes (203 of 213 targets in the 15 shown) and a model reranker raises passage recall from 0.831 to
0.922 at a fraction of a cent per question; model-written node summaries and a local cross-encoder were measured and
rejected. Tool design mattered most on unseen questions, where many fixes extended the typed query language. A higher
reasoning budget answered six more of fifteen hard questions at three times the cost; free SQL and a separate judge of
the evidence helped on failures without reaching significance (\cref{ch07}).

\section{Contributions Revisited}\label{sec:ch09-contributions}

\begin{enumerate}[nosep]
\item \textbf{A domain-neutral construction pipeline for tables and text} in which a language model proposes and code
decides: a declarative plan for the tables, a text schema, verified extraction, a mention index, and identity
resolution that asks a model only about near misses and accepts its answer only with verified quotations.
\item \textbf{The anchor graph}: a representation that uses the graph as an index of identity and evidence, with
identity as reversible, evidenced edges, a witness rule against leaked facts, and a navigation contract of five walks.
\item \textbf{An evaluation of a graph as an index}: ten criteria with hard gates fixed before measurement, judged item
by item and computed offline, with no model in the measurement loop.
\item \textbf{A measured design history}: four representations compared on the same data, showing why a fact-centric
graph underperforms for this purpose.
\item \textbf{A tool-using query engine with typed structured queries}, code-checked citations and outcomes, a
deterministic screen and a read-only layer, evaluated on blind question sets with both scores and paired tests.
\item \textbf{A low-cost, reproducible evaluation practice}: pre-registered decision rules, paired exact tests, Wilson
intervals, and every number traced to its run, its record and its commit.
\end{enumerate}

\section{Limitations}\label{sec:ch09-limitations}

\begin{itemize}[nosep]
\item The reference answers and the judgements by meaning come from one model family, with no human agreement figure.
\item The corpora are small and share one form; scale and other forms are untested.
\item Three of the four datasets were used during development; one dataset was excluded after a measurement, under a
criterion fixed in advance.
\item The headline accuracy predates the final engine's last changes, and the final engine was measured on answers
only on small sets; the outcome labels were measured on one dataset (12 questions).
\item The graph's contribution to answer accuracy was not isolated by an ablation.
\item Mention identity is per document, which fails on documents that describe several events.
\item Stored facts cover a minority of what the text states; they are kept but not relied on.
\end{itemize}

\section{Future Work}\label{sec:ch09-future}

\paragraph{Measure what is missing.} The most informative experiments are the cheapest: the engine without its graph
tools on the benchmark questions, to isolate the graph's contribution; the final engine on a fresh blind set; the
outcome set on all four datasets (48 questions not yet run); and a human check of a random sample of verdicts, to
give the judge an agreement figure.

\paragraph{Improve the engine where it fails.} The free-SQL tool is the most promising upgrade and needs a measurement
on fresh questions before adoption. The dominant remaining error, literal reading, responded only to a separate judge
of the evidence; making that judgement part of the engine, at a measured cost, is the next step for accuracy.

\paragraph{Extend the builder to other forms.} Identity within sections rather than documents, for reports that
describe several events; a hierarchy of document, section and chunk for long texts; abbreviations defined in the text
as identity evidence; and large tables kept in a database with one graph node per key. Public benchmarks that combine
tables and text, such as HybridQA~\cite{chen2020hybridqa} and TAT-QA~\cite{zhu2021tatqa}, would test the system beyond
corpora written or selected for this thesis.

\section{Closing Remarks}\label{sec:ch09-closing}

Language models make it possible to build a knowledge graph from data no one has modelled in advance, and to ask it
questions in plain language. This thesis argues, with measurements, that the most reliable use of a model in such a
system is to let it propose and let code decide; and that the graph itself is most useful not as a store of everything
the text says, but as an exact index of what the data is about and where the evidence for it lies --- with the text
read, and the counting done by a database, at the moment a question is asked.
